\section{Powering device}

 This chapter explains the steps and considerations for powering the board with a 6V battery which is supplied with the set and other ones.

\subsection{Power System Overview}
The onboard voltage regulator (MPM3610 DC-DC) steps down the input voltage on the VIN pin to 3.3V, which powers the board. The acceptable input voltage range \ref{fig:arduinopinconfig}for the VIN pin is \textbf{4.5V to 21V} .

\subsection{Supplying Power via the VIN Pin}
\begin{itemize}
    \item The \textbf{VIN pin} accepts an input voltage of 4.5V to 21V. \\
    \textbf{VIN:}  
    This pin can be used to power the board with a DC voltage source. If the power is fed through this pin, the USB power source is disconnected. This pin is an \textbf{INPUT}. Respect the voltage limits to ensure the proper functionality of the board .
    
    \item A 6V battery can be connected directly to the \textbf{VIN} pin (refer to Figure~\ref{fig:arduinopinconfig}) using a toggle switch for the positive terminal and the \textbf{GND} pin for the negative terminal. The onboard voltage regulator will convert the 6V to 3.3V to power the board .
\end{itemize}



\begin{enumerate}

    \item \textbf{Connect to VIN Pin:} Attach the positive terminal of the battery to the VIN pin thought  toggle switch and the negative terminal to the GND.\ref{fig:arduinopinconfig}
  
    \item \textbf{Test the System:} Power the board and observe the onboard LED indicators to ensure proper functionality.
\end{enumerate}






\vspace{10pt} % Optional spacing for better readability

\noindent\textbf{Source:} \href{https://store.arduino.cc/collections/boards-modules/products/arduino-nano-33-ble?_pos=19&_fid=0820327f9&_ss=c}{Arduino Nano 33 BLE Official Store Page}.
\newpage
\begin{figure}[h]
	\centering
	\includegraphics[width=1\textwidth]{Images/circuit} % Adjust the path and width as needed
	\caption{Battery Circuit}
	\label{fig:Battery_Circuit}
\end{figure}

\vspace{5pt} % Optional spacing
\noindent
\textbf{Component Description:}  
\vspace{5pt} % Optional spacing
\noindent



\begin{itemize}
    \item \textbf{G1:} 
    Battery pack (6V, 2x CR2032 cells), supplying power to the circuit.

    \item \textbf{S1:} 
    GOOBAY 10013 toggle switch:
    \begin{itemize}
        \item Connected to the positive terminal of the battery.
        \item Controls the power supply to the device (Arduino Nano 33 BLE Sense).
        \item Has two positions:
        \begin{itemize}
            \item \textbf{Open:} Disconnects the circuit, effectively cutting off the power supply.
            \item \textbf{Close:} Completes the circuit, enabling power to flow to the Arduino Nano 33 BLE Sense.
        \end{itemize}
    \end{itemize}

    \item \textbf{GND:} 
    Ground pin of the Arduino Nano 33 BLE Sense, connected to the negative terminal of the battery.\ref{fig:arduinopinconfig}

    \item \textbf{VIN:} 
    Voltage input pin of the Arduino Nano 33 BLE Sense, where the positive terminal of the battery is connected to supply power.\ref{fig:arduinopinconfig}
\end{itemize}

\begin{figure}[h]
	\centering
	\includegraphics[width=1\textwidth]{Images/circuitboard} 
	\caption{Battery Circuit and Board}
	\label{fig:BatteryCircuitBoard}
\end{figure}

\subsection{Battery Connected to the VIN Pin Alongside a USB Connection}

When a battery is connected to the VIN pin alongside a USB connection, the Arduino Nano 33 BLE Sense's power system ensures safe operation through diode isolation and voltage regulation.

The onboard diodes (\textbf{D2} and \textbf{D1}) play critical roles in this configuration:
\begin{itemize}
    \item If both USB and VIN are connected, the \textbf{VIN pin} takes priority, isolating the USB power supply. The diodes \textbf{D2} and \textbf{D1} prevent backflow from the VIN pin to the USB power line, protecting the USB source and board components.
    \item The voltage regulator (\textbf{MPM3610}) steps down the input voltage from VIN (4.5V–21V) or USB (5V via D2) to a stable 3.3V, ensuring safe operation of the board.
    \item If only USB is connected, power flows through \textbf{D2} to the VIN pin, enabling the MPM3610 regulator to function.
    \item If only a battery is connected to VIN, the regulator directly steps down the battery voltage to power the board.
\end{itemize}

The priority of power sources ensures that the board is safely powered, regardless of whether one or both sources are connected. The voltage regulator (\textbf{MPM3610}) consistently provides a stable 3.3V to the board under all conditions.


\begin{figure}[h]
	\centering
	\includegraphics[width=1\textwidth]{Images/schematic} % Adjust the path and width as needed
	\caption{Schematic}
	\label{fig:Schematic}
\end{figure}


Refer to \href{https://docs.arduino.cc/resources/schematics/ABX00031-schematics.pdf}{this schematic} for a detailed diagram showing the interaction between USB, VIN, diodes \textbf{D2} and \textbf{D1}, and the voltage regulator.


\section{Battery modification(optional)}

If you plan to use a battery with a voltage lower than 4.5V, such as a 3.3V battery, you will need to make some modifications to the board for compatibility. This ensures that the Arduino Nano 33 BLE Sense Light operates safely and effectively.

\subsection{ Using a 3.3V Battery}
When supplying a 3.3V battery, the voltage regulator on the board is bypassed. Follow these steps:

\begin{enumerate}
    \item \textbf{Direct Connection to 3.3V Pin:}
    \begin{itemize}
        \item Connect the positive\ref{fig:arduinopinconfig} terminal of the 3.3V battery to the \textbf{3.3V pin}.
        \item Connect the negative\ref{fig:arduinopinconfig} terminal to the \textbf{GND pin}.
    \end{itemize}
    
    \item \textbf{Cut the VUSB Jumper:}
    \begin{itemize}
        \item Locate the \textbf{VUSB jumper trace} on the back of the board.
        \item Use a sharp knife or a small precision tool to carefully cut the trace. Cutting this jumper ensures there is no conflict between the USB power supply and the battery.
        \item Verify the modification using a multimeter to confirm that there is no continuity between the VUSB pads.
    \end{itemize}
    \textbf{Important Notes:}
    \begin{itemize}
        \item After cutting the jumper, the USB connection will only work for data transfer (e.g., uploading code). It will no longer provide power to the board.
        \item The board must always be powered by an external battery or another suitable power source when the VUSB jumper is cut.
       \item For detailed pin configuration and additional specifications, refer to the \href{https://docs.arduino.cc/resources/pinouts/ABX00031-full-pinout.pdf}{official Arduino Nano 33 BLE pinout documentation}.
    \end{itemize}
    
 
\end{enumerate}



\begin{figure}[htbp]
	\centering
	\begin{tikzpicture}[
		pin/.style={rectangle, draw, text centered, minimum width=2.5cm, minimum height=0.8cm, font=\small},
		board/.style={draw, thick, fill=blue!20, minimum width=3cm, minimum height=15cm},
		textstyle/.style={font=\small\bfseries},
		pinpoint/.style={circle, fill=black, inner sep=0.5mm},
		every node/.style={align=center}
		]
		
		% Board
		\node[board] (board) at (0, 0) {};
		
		% USB Connector and Processor Labels
		\node[textstyle] at (0, 7) {USB};
		\node[textstyle] at (0, -7) {Processor};
		
		% Left Pins (Power, Analog, and Communication)
		\node[pin,fill=green!10] (SCK) at (-3.5, 7) {SCK};
		\node[pin,fill=red!30] (3V3) at (-3.5, 6) {+3.3V};
		\node[pin,fill=orange!30] (AREF) at (-3.5, 5) {AREF};
		\node[pin,fill=orange!30] (A0) at (-3.5, 4) {A0};
		\node[pin,fill=orange!30] (A1) at (-3.5, 3) {A1};
		\node[pin,fill=orange!30] (A2) at (-3.5, 2) {A2};
		\node[pin,fill=orange!30] (A3) at (-3.5, 1) {A3};
		\node[pin,fill=orange!30] (A4) at (-3.5, 0) {A4};
		\node[pin,fill=orange!30] (A5) at (-3.5, -1) {A5};
		\node[pin,fill=orange!30] (A6) at (-3.5, -2) {A6};
		\node[pin,fill=orange!30] (A7) at (-3.5, -3) {A7};
		\node[pin,fill=red!30] (5V) at (-3.5, -4) {+5V};
		\node[pin, fill=gray, text=white] (RESETL) at (-3.5, -5) {RESET};
		\node[pin, fill=black, text=white] (GNDL) at (-3.5, -6) {GND};
		\node[pin,fill=red!30] (VIN) at (-3.5, -7) {VIN IN};
		
		% Right Pins (Digital, SPI, and UART)
		\node[pin,fill=green!10] (D12) at (3.5, 7) {D12};
		\node[pin,fill=green!10] (D11) at (3.5, 6) {D11};
		\node[pin,fill=green!10] (D10) at (3.5, 5) {D10};
		\node[pin,fill=green!10] (D9) at (3.5, 4) {D9};
		\node[pin,fill=green!10] (D8) at (3.5, 3) {D8};
		\node[pin,fill=green!10] (D7) at (3.5, 2) {D7};
		\node[pin,fill=green!10] (D6) at (3.5, 1) {D6};
		\node[pin,fill=green!10] (D5) at (3.5, 0) {D5};
		\node[pin,fill=green!10] (D4) at (3.5, -1) {D4};
		\node[pin,fill=green!10] (D3) at (3.5, -2) {D3};
		\node[pin,fill=green!10] (D2) at (3.5, -3) {D2};
		\node[pin, fill=black, text=white] (GNDR) at (3.5, -4) {GND};
		\node[pin, fill=gray, text=white] (RESETR) at (3.5, -5) {RESET};
		\node[pin,fill=green!10] (RX) at (3.5, -6) {D1/RX};
		\node[pin,fill=green!10] (TX) at (3.5, -7) {D0/TX};
		
		% Connections (Fixed vertical alignment for corresponding pins)
		% Left Pins to Board
		\foreach \i/\y in {SCK/7, 3V3/6, AREF/5, A0/4, A1/3, A2/2, A3/1, A4/0, A5/-1, A6/-2, A7/-3, 5V/-4, RESETL/-5, GNDL/-6, VIN/-7} {
			\draw (\i.east) -- (-1.5, \y) node[pinpoint] {};
		}
		
		% Right Pins to Board
		\foreach \i/\y in {D12/7, D11/6, D10/5, D9/4, D8/3, D7/2, D6/1, D5/0, D4/-1, D3/-2, D2/-3, GNDR/-4, RESETR/-5, RX/-6, TX/-7} {
			\draw (\i.west) -- (1.5, \y) node[pinpoint] {};
		}
		
		% Add "Top View" Label
		\node[textstyle] at (0, -8) {\textbf{Top View}};
		
		% Legend
		\node[textstyle] at (0, -9) {%
			\begin{tikzpicture}
				\node[pin, fill=red!30] (power) at (0, 0) {Power Pins};
				\node[right=0.3cm of power, pin, fill=orange!30] (analog) {Analog Pins};
				\node[right=0.3cm of analog, pin, fill=green!10] (digital) {Digital Pins};
				\node[right=0.3cm of digital, pin, fill=gray, text=white] (reset) {Reset Pins};
				\node[right=0.3cm of reset, pin, fill=black, text=white] (ground) {Ground Pins};
			\end{tikzpicture}%
		};
		
	\end{tikzpicture}
	\caption{Arduino Nano 33 BLE Pins}
	\label{fig:arduinopinconfig}
\end{figure}


